% Standalone candidate-detector comparison table for the end-sem slide.
% Compile with pdflatex (or Overleaf), screenshot/export the result and
% paste into the slide as an image.
\documentclass[border=10pt]{standalone}
\usepackage{booktabs}
\usepackage{array}
\begin{document}
\small
\setlength{\tabcolsep}{6pt}
\renewcommand{\arraystretch}{1.15}
\begin{tabular}{lccccl}
\toprule
\textbf{Model} & \textbf{Params} & \textbf{Speed} & \textbf{Small-object} & \textbf{Tooling} & \textbf{Key feature} \\
               & \textbf{(M)}    &                & \textbf{detection}    & \textbf{maturity} &                       \\
\midrule
YOLOv9              & $\sim$25         & High      & Medium    & High        & PGI + GELAN backbone \\
\textbf{YOLO11-nano}& \textbf{$\sim$2.6}& \textbf{Very high} & \textbf{High} & \textbf{Very high} & \textbf{C2PSA attention block} \\
RT-DETR             & $\sim$32         & Medium    & High      & Medium      & Transformer, no NMS \\
YOLO26              & $\sim$3          & Very high & Very high & Low (new)   & End-to-end CPU inference \\
NSSRD               & Hardware-tuned   & Medium    & Very high & Low         & NAS backbone, coarse-to-fine \\
ED-Net              & $\sim$5          & High      & Very high & Low         & Dedicated XSmall head \\
\bottomrule
\end{tabular}
\end{document}
